Retrieval-augmented generation depends on finding useful evidence. The
size of the passages available for retrieval is often fixed before a
question is seen, even though different questions may need different
amounts of context. Short passages can isolate a detail but lose its
surroundings; longer passages can preserve context while introducing
irrelevant material. This thesis investigates whether adapting retrieval
granularity, the size of the retrieved chunks, improves evidence recovery
within a known source document.

The study uses QASPER scientific papers, five chunk sizes from 10 to 160
GPT-2 tokens, and fixed dense embeddings. The question-level experiments
use 2,245 training questions and 924 validation questions, with five
chunks retrieved from the associated paper. Methods progress from fixed
and mixed sizes through question-embedding and Qwen routers to similarity
trees, out-of-fold feature fusion, expected-regret learning, and local
routing. Evaluation measures lexical overlap with annotated evidence
using joined evidence-token F1.

Fixed 40-token retrieval achieves mean F1 of 0.3159. An offline oracle
that uses the reference evidence to choose the best of the five saved
results per question reaches 0.3821. The strongest learned question-level
method, expected-regret routing, reaches 0.3075, remaining below the fixed
baseline. Local routing reaches 0.3053, compared with 0.3075 for a
fixed-40 baseline using the same selected trees. Adding a frozen Qwen
context-need feature gives 0.3045, with no demonstrated improvement.

The results show that better prediction of evidence length does not
reliably translate into better retrieval. Learning from retrieval regret
produces the strongest tested global router, but the study finds no
demonstrated advantage over the corresponding fixed-size baselines.
The findings concern evidence retrieval; generated-answer quality remains
unevaluated. Repeated use of validation data during development also
limits claims about performance on unseen data.
